\section{Discussion}
\label{sec:discussion}

\paragraph{The generalist--Aristotle gap as a capability measurement.}
The gap between the best generalist ($22/60$, Gemini single-shot)
and Aristotle ($53/60$) quantifies the capability gap this benchmark
is designed to measure. Agentic mode narrows the prompt-sensitivity gap but does
not close the closure gap: Gemini agentic reaches $19/60$, still
far below $53/60$. The remaining distance is not a prompt problem.
Cross-function targets (T3) require composing definitions from
multiple library modules in a single proof step. Generalists
trained on Mathlib do not learn this composition because Mathlib
contains no cross-family CS inequalities. The T3 hint ablation
(Appendix~\ref{app:hint-ablation-table}) confirms that supplying
both the substitution and the lemma chain closes the gap
($3$--$5/6$ at \textsc{Sub+Chain}), but the drafter cannot
produce those hints on its own.

\paragraph{Current provers are memorisation engines, not reasoners.}
Goedel-Prover-V2~\cite{goedelprover2025}, trained on Mathlib,
responds to \texttt{etabetting\_nonneg} with \emph{``the Lean~4
code provided does not define etaBetting \ldots\ we cannot proceed
with a meaningful proof.''} This is not a failure of tactic
knowledge. It is a failure to read a new definition and reason
about it. Any domain-specific Lean library outside Mathlib will
surface this gap: the prover has seen the tactics but never the
definitions. Retrieval-augmented proving, fine-tuning on
domain libraries, or longer-context in-context learning are
candidate next steps. The benchmark provides the measurement
tool. The T1 tier ($34$ single-function targets against custom
definitions) is designed to isolate this axis.

\paragraph{Refutation as a safety signal.}
Aristotle surfaces three false-as-stated targets that every other
solver accepted. Two are definitional errors (an off-by-one in a
stopping-rule admissibility condition and a missing monotonicity
hypothesis); one is a quantitative counterexample. In all three
cases, the generalist drafters produce confident, syntactically
valid Lean proofs that the kernel rejects only because the
statement is false. A deployment pipeline that trusts the drafter
and skips kernel verification would ship a broken guarantee. The
refutation track measures a capability axis absent from pass@$N$
benchmarks: whether a solver can identify that a statement is
wrong, not just fail to prove it.

\paragraph{Limitations.}
The benchmark covers one subfield (anytime-valid CS) with $60$
targets. Closure rates may not transfer to other Lean libraries
with different import structures or proof shapes. Aristotle is a
closed-source system. Its results are reproducible via session IDs
but not by re-running an open model. The $60$-target scale limits
statistical power on per-tier comparisons. Wilson 95\% CIs
and contamination notes are in Appendix~\ref{tab:wilson-ci}.

\paragraph{Broader impact.}
Anytime-valid confidence sequences are deployed in production A/B
testing and clinical trial monitoring systems that make decisions
affecting millions of users. A silent coverage-guarantee violation
in these systems can lead to incorrect treatment decisions or
product launches based on false statistical significance. Formally
verified implementations remove this failure class. This benchmark
measures how close current language-model provers are to automating
that verification, and the refutation track provides an early
warning for the complementary risk: solvers that produce
plausible-looking proofs of false statements.
